% This document was generated by an LLM.

\documentclass{workbook}

\usepackage{amsthm}
\usepackage{booktabs}

\newtheoremstyle{step}{\topsep}{\topsep}{\normalfont}{0pt}{\bfseries}{.}{0.5em}{}
\theoremstyle{step}
\newtheorem*{question}{Question}
\newtheorem*{hint}{Hint}

\newcommand{\dd}{\,\mathrm{d}}

\title{Stirling's Formula\\A Guided Derivation}
\author{}
\date{}

\begin{document}

\maketitle

\section*{How to use this document}

The goal is to derive
\[
  n! \sim \sqrt{2\pi n}\left(\frac{n}{e}\right)^{n},
  \qquad\text{that is,}\qquad
  \lim_{n\to\infty}\frac{n!}{\sqrt{2\pi n}\,(n/e)^{n}} = 1 .
\]
The derivation is split into eight steps. Each step states a question and
gives a hint. Attempt every step before turning to the solutions in
Section~\ref{sec:solutions}. The solutions are collected separately so that
the result of a step does not give away how to reach it.

Steps 1--5 and 8 use only single-variable calculus: integration by parts,
derivatives, Taylor expansion, and substitution. Step~6 uses a double
integral in polar coordinates. Step~7 is optional; it makes the limit argument
rigorous using elementary inequalities.

\section{The plan}

The individual steps are routine. The difficulty is knowing which steps to
take. One idea determines all of them.

\medskip
\noindent\textbf{Laplace's method.} If a positive function has a single
sharp peak, its integral is approximately the height of the peak times its
width. Near a maximum, the logarithm of a smooth function is approximately a
downward parabola, so the function itself is approximately a bell curve. The
integral of a bell curve is known exactly.

\medskip
\noindent This suggests the following road map.
\begin{enumerate}
  \item Express $n!$ as an integral.
  \item Find where the integrand peaks and how tall the peak is.
  \item Measure the curvature of the logarithm at the peak.
  \item Rescale the variable so the peak has a shape independent of $n$.
  \item Show that the rescaled integrand converges to a bell curve.
  \item Evaluate the integral of the limiting bell curve.
  \item (Optional) Justify taking the limit inside the integral.
  \item Assemble the factors.
\end{enumerate}
Each of the three factors in Stirling's formula appears at a definite point
on this road map: $(n/e)^n$ in Step~2, $\sqrt{n}$ in Step~4, and
$\sqrt{2\pi}$ in Step~6.

\section{The steps}

\subsection*{Step 1: Turn the product into an integral}

\begin{question}
The product $n! = 1\cdot 2\cdots n$ is discrete, and calculus does not apply
to it directly. Find an integral $I_n$ satisfying $I_n = n\,I_{n-1}$ and
$I_0 = 1$. Any such sequence equals $n!$.
\end{question}

\begin{hint}
Integration by parts lowers a power: differentiating $x^n$ produces the
factor $n$ and leaves $x^{n-1}$. Pair $x^n$ with a function that keeps its
form under integration, so the new integral has the same shape as the old
one. Choose the limits of integration so the boundary terms vanish.
\end{hint}

\subsection*{Step 2: Locate the peak}

\begin{question}
Sketch the integrand from Step~1. One factor grows and the other decays, so
the product rises and then falls. Where is the maximum, and what is its
value?
\end{question}

\begin{hint}
Products are awkward to differentiate. Write the integrand as $e^{f(x)}$ and
maximize $f$ instead.
\end{hint}

\subsection*{Step 3: Measure the curvature at the peak}

\begin{question}
How quickly does $f$ fall off on either side of its maximum?
\end{question}

\begin{hint}
Expand $f$ in a Taylor series around the maximum. The first-order term
vanishes at a maximum, so the second-order term is the first one that
carries information. What does $e^{f(x)}$ look like when $f$ is replaced by
this quadratic approximation?
\end{hint}

\subsection*{Step 4: Rescale so the peak has a fixed shape}

\begin{question}
The bell curve from Step~3 becomes wider as $n$ grows. Find a change of
variables after which its shape no longer depends on $n$.
\end{question}

\begin{hint}
Substitute $x = n + s\,t$ with a scale factor $s$ to be determined. Choose $s$
so that the quadratic exponent becomes exactly $-t^2/2$. Remember that the
substitution also changes $\dd x$.
\end{hint}

\subsection*{Step 5: Substitute exactly and examine the error}

\begin{question}
Steps 3 and 4 used an approximation. Perform the substitution from Step~4 in
$f$ exactly, without truncating anything. Write
$n!/\bigl(\sqrt{n}\,(n/e)^n\bigr)$ as an integral over $t$. What does the
integrand converge to as $n\to\infty$, and what happens to the limits of
integration?
\end{question}

\begin{hint}
Factor $n$ out of the logarithm:
$\ln(n + \sqrt{n}\,t) = \ln n + \ln\bigl(1 + t/\sqrt{n}\bigr)$. Then use
$\ln(1+y) = y - \tfrac{y^2}{2} + \tfrac{y^3}{3} - \cdots$ and track the power
of $n$ attached to each term.
\end{hint}

\subsection*{Step 6: Evaluate the limiting integral}

\begin{question}
Step 5 suggests that
\[
  \frac{n!}{\sqrt{n}\,(n/e)^n} \;\longrightarrow\; G
  := \int_{-\infty}^{\infty} e^{-t^2/2}\dd t .
\]
Compute $G$.
\end{question}

\begin{hint}
The function $e^{-t^2/2}$ has no elementary antiderivative, so compute $G^2$
instead. Write it as a double integral over the plane in variables $t$ and
$u$. The exponent depends only on $t^2 + u^2$, so switch to polar
coordinates. The area element $r\dd r\dd\theta$ supplies the factor $r$ that
makes the radial integral elementary.
\end{hint}

\subsection*{Step 7 (optional): Justify passing to the limit}

\begin{question}
Step 5 shows that the integrand converges for each fixed $t$. This does not
by itself imply that the integrals converge, because mass could escape to
infinity in the tails. Find a single integrable function, independent of
$n$, that bounds the integrand from Step~5 for all $n \ge 1$. Then use it to
prove that the integrals converge to $G$.
\end{question}

\begin{hint}
Prove the two inequalities
\begin{align*}
  \ln(1+y) &\le y - \frac{y^2}{2}      && \text{for } -1 < y \le 0, \\
  \ln(1+y) &\le y - \frac{y^2}{2(1+y)} && \text{for } y \ge 0 .
\end{align*}
For each, define the difference of the two sides, check that it vanishes at
$y = 0$, and determine the sign of its derivative. Apply them with
$y = t/\sqrt{n}$ and multiply by $n$. To finish, split the integral into a
bounded interval $[-A, A]$ and the two tails.
\end{hint}

\subsection*{Step 8: Assemble}

\begin{question}
Combine Steps 5 and 6 to obtain Stirling's formula. Then compute the ratio
\[
  \frac{\sqrt{2\pi n}\,(n/e)^n}{n!}
\]
for $n = 1, 5, 10, 20$, and compare the relative error with $1/(12n)$.
\end{question}

\begin{hint}
If $a_n/b_n \to G$ and $G \ne 0$, then $a_n / (G\,b_n) \to 1$.
\end{hint}

\section{Solutions}\label{sec:solutions}

\subsection*{Step 1}

Let $I_n = \int_0^\infty x^n e^{-x}\dd x$. Then $I_0 = 1$, and integration by
parts with $u = x^n$, $\dd v = e^{-x}\dd x$ gives
\[
  I_n = \Bigl[-x^n e^{-x}\Bigr]_0^\infty + n\int_0^\infty x^{n-1}e^{-x}\dd x
      = n\,I_{n-1},
\]
since $x^n e^{-x}\to 0$ as $x\to\infty$ and vanishes at $x=0$ for $n\ge 1$.
Hence
\[
  n! = \int_0^\infty x^n e^{-x}\dd x .
\]

\subsection*{Step 2}

Write $x^n e^{-x} = e^{f(x)}$ with $f(x) = n\ln x - x$. Then
$f'(x) = n/x - 1$, which vanishes at $x = n$. The peak height is
\[
  e^{f(n)} = e^{n\ln n - n} = \left(\frac{n}{e}\right)^{n}.
\]
This is the first factor of Stirling's formula.

\subsection*{Step 3}

Since $f''(x) = -n/x^2$, we have $f''(n) = -1/n$, and
\[
  f(x) \approx f(n) - \frac{(x-n)^2}{2n},
  \qquad\text{so}\qquad
  x^n e^{-x} \approx \left(\frac{n}{e}\right)^{n} e^{-(x-n)^2/(2n)} .
\]
Near its maximum the integrand is approximately a bell curve centred at $n$.

\subsection*{Step 4}

With $x = n + s\,t$, the exponent becomes $-s^2 t^2/(2n)$. Setting
$s^2/(2n) = 1/2$ gives $s = \sqrt{n}$. Thus
\[
  x = n + \sqrt{n}\,t, \qquad \dd x = \sqrt{n}\dd t .
\]
The factor $\sqrt{n}$ is the width of the peak. It is the second factor of
Stirling's formula.

\subsection*{Step 5}

The exact substitution gives
\[
  f(n + \sqrt{n}\,t) = n\ln n - n + g_n(t),
  \qquad
  g_n(t) = n\ln\!\left(1 + \frac{t}{\sqrt{n}}\right) - \sqrt{n}\,t .
\]
The range $0 < x < \infty$ corresponds to $-\sqrt{n} < t < \infty$, so
\[
  \frac{n!}{\sqrt{n}\,(n/e)^n} = \int_{-\sqrt{n}}^{\infty} e^{g_n(t)}\dd t .
\]
This identity is exact. Expanding the logarithm,
\[
  g_n(t) = -\frac{t^2}{2} + \frac{t^3}{3\sqrt{n}} - \frac{t^4}{4n} + \cdots
  \;\longrightarrow\; -\frac{t^2}{2}
  \qquad (n\to\infty)
\]
for each fixed $t$. The linear terms cancel because the expansion is taken at
the maximum, and every term beyond the quadratic one carries a negative power
of $n$. The lower limit $-\sqrt{n}$ tends to $-\infty$.

\subsection*{Step 6}

\[
  G^2 = \int_{-\infty}^{\infty}\!\int_{-\infty}^{\infty}
        e^{-(t^2+u^2)/2}\dd t\dd u
      = \int_0^{2\pi}\!\int_0^{\infty} e^{-r^2/2}\, r\dd r\dd\theta
      = 2\pi \Bigl[-e^{-r^2/2}\Bigr]_0^{\infty}
      = 2\pi .
\]
Since $G > 0$, we get $G = \sqrt{2\pi}$. This is the third factor of
Stirling's formula, and the only step that uses a double integral.

\subsection*{Step 7}

\emph{First inequality.} Let $k(y) = y - y^2/2 - \ln(1+y)$ on $(-1, 0]$. Then
$k(0) = 0$ and
\[
  k'(y) = 1 - y - \frac{1}{1+y} = -\frac{y^2}{1+y} \le 0 .
\]
So $k$ is non-increasing, and $k(y) \ge k(0) = 0$ for $y \le 0$.

\emph{Second inequality.} Let
$h(y) = y - \dfrac{y^2}{2(1+y)} - \ln(1+y)$ on $[0,\infty)$. Then $h(0) = 0$
and a direct computation gives
\[
  h'(y) = \frac{y^2}{2(1+y)^2} \ge 0 ,
\]
so $h(y) \ge 0$ for $y \ge 0$.

\emph{Bound on $g_n$.} Note that
$g_n(t) = n\bigl[\ln(1+y) - y\bigr]$ with $y = t/\sqrt{n}$. For $t\le 0$ the
first inequality gives $g_n(t) \le -t^2/2$. For $t \ge 0$ the second gives
\[
  g_n(t) \le -\frac{t^2}{2(1 + t/\sqrt{n})} \le -\frac{t^2}{2(1+t)}
  \qquad (n \ge 1).
\]
Define
\[
  H(t) =
  \begin{cases}
    e^{-t^2/2}, & t \le 0,\\[0.3em]
    e^{-t^2/(2(1+t))}, & t \ge 0.
  \end{cases}
\]
Then $e^{g_n(t)} \le H(t)$ for all $n\ge 1$ and all $t > -\sqrt{n}$, and $H$
is integrable over $\mathbb{R}$: for $t \ge 1$ we have
$t^2/(2(1+t)) \ge t/4$, so $H(t) \le e^{-t/4}$.

\emph{Convergence.} Let $\varepsilon > 0$. Choose $A$ so large that
$\int_{|t|>A} H < \varepsilon$ and $\int_{|t|>A} e^{-t^2/2} < \varepsilon$.
Both tail contributions are then less than $\varepsilon$ for every $n$. On
$[-A, A]$, Taylor's theorem with remainder shows that
$g_n(t) \to -t^2/2$ uniformly, because the error is bounded by a constant
times $A^3/\sqrt{n}$ once $\sqrt{n} > 2A$. Hence the integrals over
$[-A, A]$ converge, and for all sufficiently large $n$,
\[
  \left| \int_{-\sqrt{n}}^{\infty} e^{g_n(t)}\dd t - \sqrt{2\pi}\right|
  < 3\varepsilon .
\]
This argument is a special case of the dominated convergence theorem.

\subsection*{Step 8}

Combining Steps 5, 6 and 7,
\[
  \frac{n!}{\sqrt{n}\,(n/e)^n} \longrightarrow \sqrt{2\pi},
  \qquad\text{hence}\qquad
  n! \sim \sqrt{2\pi n}\left(\frac{n}{e}\right)^{n}.
\]

\begin{center}
\begin{tabular}{rrcc}
  \toprule
  $n$ & $n!$ & $\sqrt{2\pi n}\,(n/e)^n / n!$ & $1/(12n)$ \\
  \midrule
   1 & $1$                    & $0.92214$ & $0.08333$ \\
   5 & $120$                  & $0.98349$ & $0.01667$ \\
  10 & $3\,628\,800$          & $0.99170$ & $0.00833$ \\
  20 & $\approx 2.4329\times 10^{18}$ & $0.99584$ & $0.00417$ \\
  \bottomrule
\end{tabular}
\end{center}

\noindent The relative error $1 - \sqrt{2\pi n}\,(n/e)^n/n!$ is close to
$1/(12n)$, which is the first correction term in the full Stirling series
$n! = \sqrt{2\pi n}\,(n/e)^n\bigl(1 + \tfrac{1}{12n} + O(n^{-2})\bigr)$.

\section{Remarks}

\paragraph{The general method.} Steps 2--6 apply to any integral of the form
$\int e^{n\phi(x)}\dd x$ where $\phi$ has a single interior maximum $x_0$
with $\phi''(x_0) < 0$:
\[
  \int e^{n\phi(x)}\dd x \sim e^{n\phi(x_0)}\sqrt{\frac{2\pi}{n\,|\phi''(x_0)|}} .
\]
For $n!$, substitute $x = n y$ to obtain
$n! = n^{n+1}\int_0^\infty e^{n(\ln y - y)}\dd y$, with $\phi(y) = \ln y - y$,
$y_0 = 1$, $\phi''(1) = -1$.

\paragraph{An alternative route.} Comparing
$\ln n! = \sum_{k=1}^{n}\ln k$ with $\int_1^n \ln x\dd x$ and applying the
trapezoid rule gives
\[
  \ln n! = n\ln n - n + \tfrac{1}{2}\ln n + C + o(1)
\]
for some constant $C$. This route is elementary, but it does not determine
$C$. Showing that $e^{C} = \sqrt{2\pi}$ requires an independent input such as
the Wallis product. The integral method above produces the constant
automatically.

\end{document}
